\SourceLeaf{99}%
introduce a new function, he had chosen the one above in preference to his li~$x = \int \frac{dx}{\log x}$. A
single-valued function is invariably to be regarded as incomparably more classical and simpler than
a multivalued one, especially since $\log x$ is itself already a multivalued function. $\mathord{.}\,\mathord{.}\,\mathord{.}\,\mathord{.}\,\mathord{.}\,\mathord{.}\,\mathord{.}\,\mathord{.}\,\mathord{.}$

$\mathord{.}\,\mathord{.}\,\mathord{.}\,\mathord{.}$ Furthermore, I believe that extending the investigations to imaginary
arguments will lead to highly interesting results. For the reasons given above, however, I would
prefer to choose the function $\int \frac{e^x-1}{x}\,dx$ over
$\int \frac{dx}{\log x}$, because I suspect the former will yield more elegant results.
For example, I would very much like to know whether that function or, equivalently,
the series

\[
x + \tfrac14 xx + \tfrac{1}{18}x^3 + \text{etc.}
\]

\noindent can, for certain finite values of $x$ of the form $a+bi$, become $0$. I
cannot yet assert this with certainty, although it seems very probable to me. If there are
such values (and then there are certainly infinitely many), they will be very remarkable quantities,
and the entire series will factor into infinitely many factors of the form $(1+2ax+\beta xx)$.
---

$\mathord{.}\,\mathord{.}\,\mathord{.}\,\mathord{.}\,\mathord{.}\,\mathord{.}\,\mathord{.}\,\mathord{.}\,\mathord{.}\,\mathord{.}\,\mathord{.}\,\mathord{.}\,\mathord{.}$

\Needspace{14\baselineskip}
\begin{center}
REMARK.
\end{center}

\noindent\hspace*{\parindent}This discussion of the foundations of the theory of functions of a complex variable
is taken, as the text shows, from a letter to \textsc{Bessel} of 18\ December 1811
that has already been published in the »Briefwechsel zwischen \textsc{Gauss} und \textsc{Bessel}«
(Leipzig, 1880). It is probable, moreover, that the view
developed here by \textsc{Gauss} had been familiar to him long before he wrote the
letter in question. \textsc{Gauss} presumably pursued, for\,example,\ the investigations in the
second half of
Art.\ 12 of the »Arithmetisch-geometrischen Mittel« (Collected\ Works III p.\ 378) on the basis of these very views.
Those investigations probably belong to a much earlier period. \hfill \textsc{Fricke}.
